% from main.tex
%%%%%%%%%%  scan idx 512  |  folio 499  %%%%%%%%%%
\[
K_\bullet(X')\xrightarrow{i_*}K_\bullet(X)\xrightarrow{j^*}K_\bullet(U)\longrightarrow 0
\]
est
exacte.

Si $F'$ est un $\mathcal O_{X'}$-Module cohérent, $i_*(F')$ est un $\mathcal O_X$-Module dont le support est contenu dans $X'$; par conséquent sa restriction à $U$ est nulle, et $j^*\circ i_*(\cl(F'))=0$, d'où, les éléments de la forme $\cl(F')$ engendrant $K_\bullet(X')$, $j^*\circ i_*=0$. On en déduit un homomorphisme
\[
 \rho:K_\bullet(X)/\operatorname{Im}K_\bullet(X')\longrightarrow K_\bullet(U).
\]
On va définir un homomorphisme inverse de $\rho$. Soit $G$ un $\mathcal O_U$-Module cohérent. Comme $X$ est noethérien, il existe un $\mathcal O_X$-Module cohérent $F$ dont la restriction à $U$ est isomorphe à $G$. Si $F'$ est un second $\mathcal O_X$-Module cohérent dont la restriction à $U$ est isomorphe à $G$, alors on a dans $K_\bullet(X)$ :
\[
 \cl(F)-\cl(F')\in \operatorname{Im}K_\bullet(X').
\]
Pour le voir, on considère $G$ comme un sous-faisceau de $j^*(F\oplus F')$ au moyen de l'application diagonale; $G$ se prolonge alors en un sous-faisceau $F''$ de $F\oplus F'$, et on est ramené à montrer que
\[
 \cl(F)-\cl(F'')\in \operatorname{Im}K_\bullet(X'),
\]
et l'analogue pour $F'$. Soient alors $H$ et $K$ le noyau et le conoyau de l'homomorphisme composé :
\[
F''\longrightarrow F\oplus F'\longrightarrow F ;
\]
cet homomorphisme se réduit à l'identité de $G$ sur $U$, et par suite $H$ et $K$ sont à support contenu dans $X'$. Si $I$ est l'idéal de $X'$ dans $X$, il en résulte qu'il existe des entiers $m$ et $n$ tels que $I^m\cdot H=0$, $I^n\cdot K=0$ (EGA I 9.3.4). On peut donc dévisser $H$ et $K$ en extensions de Modules
